\documentclass[11pt,a4paper]{article}

\usepackage[utf8]{inputenc}
\usepackage[T1]{fontenc}
\usepackage{lmodern}
\usepackage[margin=1.1in]{geometry}
\usepackage{graphicx}
\usepackage{booktabs}
\usepackage{float}
\usepackage{xcolor}
\usepackage{hyperref}
\usepackage{enumitem}
\usepackage{amsmath}
\usepackage{mdframed}
\usepackage{subcaption}

\hypersetup{colorlinks=true, linkcolor=blue!60!black, urlcolor=blue!60!black}

\definecolor{tuvblue}{RGB}{44,90,160}
\definecolor{lightblue}{RGB}{220,235,255}
\definecolor{lightgreen}{RGB}{220,245,220}
\definecolor{lightyellow}{RGB}{255,250,220}
\definecolor{lightred}{RGB}{255,230,230}
\definecolor{lightgray}{RGB}{245,245,245}

\newmdenv[backgroundcolor=lightblue,linecolor=tuvblue,linewidth=1.5pt,
  innerleftmargin=12pt,innerrightmargin=12pt,
  innertopmargin=8pt,innerbottommargin=8pt]{keybox}
\newmdenv[backgroundcolor=lightgreen,linecolor=green!40!black,linewidth=1pt,
  innerleftmargin=12pt,innerrightmargin=12pt,
  innertopmargin=8pt,innerbottommargin=8pt]{findingbox}
\newmdenv[backgroundcolor=lightyellow,linecolor=orange!60!black,linewidth=1pt,
  innerleftmargin=12pt,innerrightmargin=12pt,
  innertopmargin=8pt,innerbottommargin=8pt]{notebox}
\newmdenv[backgroundcolor=lightred,linecolor=red!40!black,linewidth=1pt,
  innerleftmargin=12pt,innerrightmargin=12pt,
  innertopmargin=8pt,innerbottommargin=8pt]{warningbox}
\newmdenv[backgroundcolor=lightgray,linecolor=gray!50,linewidth=1pt,
  innerleftmargin=12pt,innerrightmargin=12pt,
  innertopmargin=8pt,innerbottommargin=8pt]{codebox}

\title{\textbf{Agent Efficiency Experiment: Results Report}\\[0.4em]
\large Code Coverage on Spring PetClinic\\
7 Variants $\times$ 3 Runs --- What the Data Shows\\[0.6em]
\normalsize Tuvium Code Coverage Experiment v2}
\author{Tuvium Research}
\date{March 2026}

\begin{document}
\maketitle
\tableofcontents
\newpage

%% ============================================================
\section{What We Ran}

\subsection*{The task}

Write JUnit 5 tests for Spring PetClinic from scratch to maximize code coverage.
The agent starts with zero tests.
It finishes when it decides the coverage is satisfactory --- or when it times out.

Spring PetClinic is the reference Spring Boot application: well-documented,
heavily covered in online tutorials, present in training data for all major language models.
This matters for interpreting the results (see Section~\ref{sec:limits}).

\subsection*{The setup}

\begin{itemize}
  \item \textbf{Model}: Claude Sonnet 4.6
  \item \textbf{Task}: write tests for spring-petclinic (zero initial coverage)
  \item \textbf{Coverage target}: instruction coverage via JaCoCo
  \item \textbf{Build}: \texttt{./mvnw clean test jacoco:report}
  \item \textbf{Runs}: 20 sessions total (N=3 per variant, except \texttt{simple}: N=2)
  \item \textbf{Tracking}: every tool call logged in order via \texttt{agent-journal}
\end{itemize}

\subsection*{The seven variants}

Each variant gives the agent a different level of structured support.
They build additively from a minimal prompt up to a two-phase exploration+action structure.

\begin{table}[H]
\centering
\renewcommand{\arraystretch}{1.2}
\begin{tabular}{@{}l p{9cm}@{}}
\toprule
\textbf{Variant} & \textbf{What the agent has} \\
\midrule
\texttt{simple}
  & Minimal prompt. No stopping condition. No knowledge.
    The bare minimum needed to run the task. \\
\texttt{hardened}
  & Structured prompt with explicit stopping condition and output format.
    Still no domain knowledge. \\
\texttt{hardened+kb}
  & Hardened prompt plus a flat knowledge base: Spring test annotation imports,
    JaCoCo configuration, petclinic package structure. \\
\texttt{hardened+sae}
  & Hardened prompt plus Structured Analysis and Exploration (pre-analysis):
    a mandatory pre-analysis pass before writing any tests. \\
\texttt{hardened+skills}
  & Hardened prompt plus SkillsJars: structured, modular knowledge packages
    delivered as callable skills. \\
\texttt{hardened+skills+sae}
  & Skills plus pre-analysis together: structured knowledge delivery and
    mandatory pre-analysis. \\
\texttt{hardened+skills+sae+plan-act}
  & Two-phase design: a deep exploration phase followed by a sustained act phase.
    The most structured variant. \\
\bottomrule
\end{tabular}
\caption{Seven variants tested, ordered from least to most structured.}
\end{table}

\subsection*{What we measured}

Every tool call was classified into one of 9 semantic states
(EXPLORE, SHELL, READ\_KB, READ\_SKILL, JAR\_INSPECT, WRITE, BUILD, VERIFY, FIX).
SHELL captures shell-based searching (find, ls, grep, cat) --- unstructured exploration
distinct from EXPLORE (Read/Glob tool calls that go directly to known locations).
See the companion document \emph{Reading Agent Behavior} for the full taxonomy.

At the end of each run, a separate judge (T3 TestQualityJudge) evaluated
the produced test suite on 6 criteria: test slice selection, assertion quality,
error coverage, domain patterns, coverage target selection, and version-aware patterns.
Score range: 0.0--1.0. Runs were not scored by the same process that produced the tests
(no circular grading).

%% ============================================================
\section{The Stark Contrast}

Before the summary tables: two real runs, same model, same task.

\begin{center}
\renewcommand{\arraystretch}{1.3}
\begin{tabular}{@{}lcccc@{}}
\toprule
& \textbf{Total calls} & \textbf{BUILD calls} & \textbf{Cost} & \textbf{Duration} \\
\midrule
\textbf{hardened+skills+sae} (run 17) & 69 & 11 & \$3.49 & 26 min \\
\textbf{hardened+skills+sae+plan-act} (run 14) & 97 & \textbf{23} & \textbf{\$5.67} & 43 min \\
\bottomrule
\end{tabular}
\end{center}

Run 17 built once, verified, done.
Run 14 built 23 times --- six consecutive FIX$\to$BUILD cycles in the middle of the run.
The 63\% cost difference comes entirely from that spiral.

Both runs produced passing tests.
Both scored above 0.85 on the T3 judge.
The spiral did not improve the outcome.

\begin{figure}[H]
\centering
\includegraphics[width=\textwidth]{../figures/cost-steps-comparison.png}
\caption{Mean cost (left) and mean step count (right) per variant across all runs.
         The plan-act variant (red bar) has the highest mean cost and the highest variance ---
         it can be the cleanest or the most expensive run in the dataset depending
         on whether it hits a spiral. Error bars = standard deviation across N=3 runs.}
\end{figure}

%% ============================================================
\section{Two-Axis Results}
\label{sec:results}

Efficiency and quality are independent axes.
A variant that saves 30\% of steps but degrades quality isn't actually saving you money
on the goal of \emph{producing good tests}.

\begin{table}[H]
\centering
\renewcommand{\arraystretch}{1.3}
\begin{tabular}{@{}lrrrrrr@{}}
\toprule
\textbf{Variant} & \textbf{N} & \textbf{Mean steps} & \textbf{$\pm$} &
  \textbf{Mean cost} & \textbf{T3 quality} & \textbf{$\pm$} \\
\midrule
\texttt{hardened+skills+sae}     & 3 &  75.0 & 8.7  & \$3.39 & 0.850 & 0.017 \\
\texttt{hardened+sae}            & 3 &  80.3 & 12.1 & \$3.41 & 0.789 & 0.035 \\
\texttt{hardened+kb}             & 3 &  83.3 & 7.4  & \$3.21 & 0.847 & 0.027 \\
\texttt{hardened+skills+sae+plan-act} & 3 & 95.0 & 11.1 & \$5.11 & 0.878 & 0.092 \\
\texttt{hardened+skills}         & 3 & 101.7 & 10.1 & \$3.70 & 0.856 & 0.025 \\
\texttt{hardened}                & 3 & 103.0 & 17.8 & \$4.08 & 0.850 & 0.017 \\
\texttt{simple}                  & 2 & 109.5 & 13.4 & \$3.47 & 0.783 & 0.071 \\
\bottomrule
\end{tabular}
\caption{Two-axis comparison: efficiency (steps, cost) and quality (T3 score).
         Steps = total tool calls per run. T3 quality = mean across N runs $\pm$ std.
         Sorted by mean steps ascending (most efficient first).}
\end{table}

\begin{warningbox}
\textbf{The plan-act paradox}: \texttt{hardened+skills+sae+plan-act} ranks 4th on efficiency
(95 mean steps) but has the highest mean quality (0.878) \emph{and} the highest standard
deviation (0.092). Its best run (run 15: \$5.80, T3=0.967) and worst run (run 14:
\$5.67, 23 builds, spiral) are the extremes of the entire dataset.
High ceiling, high floor risk.
\end{warningbox}

%% ============================================================
\section{Four Key Findings}

\begin{findingbox}
\textbf{Finding 1: Prompt hardening is the biggest quality driver.}

\texttt{simple} $\to$ \texttt{hardened}: $+$0.067 quality points (0.783 $\to$ 0.850),
$-$6\% steps. This is the largest single quality jump in the table.

All other knowledge and structure variants build on top of a hardened prompt.
The hardened prompt adds a stopping condition and structured output requirements ---
no external knowledge, just discipline.
\emph{Structured prompting beats no structure, regardless of what knowledge you inject.}
\end{findingbox}

\begin{findingbox}
\textbf{Finding 2: pre-analysis drives efficiency at a quality cost.}

\texttt{hardened} $\to$ \texttt{hardened+sae}: $-$22\% steps,
but quality drops to 0.789 --- back near \texttt{simple}'s 0.783.

pre-analysis's mandatory pre-analysis pass front-loads exploration and produces a tighter
action plan. The agent is more efficient: it writes fewer exploratory tests,
makes fewer speculative BUILD attempts.
But the plan can be too tight --- it scopes out edge cases and domain-specific
behaviors that a less constrained agent would discover during exploration.
\emph{Efficiency and depth are in tension. pre-analysis resolves it toward efficiency.}
\end{findingbox}

\begin{findingbox}
\textbf{Finding 3: Skills compensate for pre-analysis's quality regression.}

\texttt{hardened+skills+sae}: $-$31\% steps vs simple, quality = 0.850 ---
matching \texttt{hardened} alone, despite being 28 steps shorter.

Skills provide structured domain knowledge: correct test patterns, annotation imports,
petclinic-specific coverage strategies. With skills, the pre-analysis plan is better-informed ---
it scopes tasks correctly because it has the right vocabulary to scope with.
The skill package functions as a quality stabilizer when combined with pre-analysis.
\emph{This is the strongest variant by the efficiency+quality tradeoff.}
\end{findingbox}

\begin{findingbox}
\textbf{Finding 4: KB is a pure efficiency play on known codebases.}

\texttt{hardened+kb}: $-$24\% steps vs \texttt{hardened}, quality = 0.847 (flat,
$\Delta$ = $-$0.003, within noise).

The knowledge base reduces JAR inspection (10.1\% $\to$ 3\% of steps roughly)
by answering the import questions the agent would otherwise search for.
But petclinic is well-known --- the model already knows its domain,
so KB cannot improve the tests, only reduce friction.
\emph{On a novel codebase, KB's quality impact should be larger.
This is the central hypothesis for the Fineract follow-up.}
\end{findingbox}

%% ============================================================
\section{The Behavioral Signature: What the Heatmaps Show}

The two-axis table shows what happened. The heatmaps show \emph{why}.

\subsection*{The JAR cluster: knowledge friction}

When the agent doesn't know which Maven jar provides a Spring test annotation,
it inspects jars in \texttt{\textasciitilde/.m2/repository} until it finds the class.
This is addressable: a knowledge base entry answering the import question
eliminates it entirely.

\begin{center}
\begin{tabular}{@{}lrrc@{}}
\toprule
\textbf{Variant} & \textbf{JAR\_INSPECT \%} & \textbf{SEARCH \%} & \textbf{Intervention} \\
\midrule
\texttt{simple}              &  9.3\% & 24.8\% & none \\
\texttt{hardened}            & 17.7\% & 38.0\% & prompt only (more thorough, no direction) \\
\texttt{hardened+kb}         & 15.0\% & 30.8\% & flat KB helps but doesn't eliminate \\
\texttt{hardened+skills}     &  1.7\% & 13.4\% & skills answer the import directly \\
\texttt{hardened+skills+sae} &  8.5\% & 17.9\% & pre-analysis reduces search overhead \\
\bottomrule
\end{tabular}
\end{center}

SEARCH = SHELL + JAR\_INSPECT (all unstructured searching). This is the headline efficiency
metric: \texttt{hardened} wastes 38\% of all steps on searching; \texttt{hardened+skills}
cuts that to 13.4\%, and direct EXPLORE (targeted file reads) doubles from 30\% to 60\%.
Skills don't just reduce JAR inspection --- they shift the agent from searching to reading.

\begin{notebox}
\textbf{The hardened paradox}: adding a hardened prompt \emph{increases} JAR inspection
from 9.3\% to 17.7\%. The hardened prompt makes the agent more thorough about finding
correct imports --- without providing the answer. Thoroughness without direction =
maximum searching.
\end{notebox}

\subsection*{The FIX/BUILD loop: rework}

When tests fail to compile or run, the agent edits (FIX) and rebuilds (BUILD).
One or two FIX$\to$BUILD cycles is normal. Six consecutive cycles is a spiral.

The heatmap for \texttt{hardened+skills+sae+plan-act} shows mutual high probability
between the BUILD and FIX rows: $P(\text{BUILD}\to\text{FIX})$ and
$P(\text{FIX}\to\text{BUILD})$ both elevated. Compare to \texttt{hardened+skills+sae}:
the FIX row is dimmer, fewer cycles, cleaner exit.

The plan-act variant's two-phase design (deep explore, then sustained act) produces
more ambitious initial test writing. When that writing hits a compilation problem,
the agent retries aggressively. The recovery is expensive precisely because the initial
attempt was so thorough.

\begin{figure}[H]
\centering
\begin{minipage}[t]{0.48\textwidth}
  \includegraphics[width=\textwidth]{../figures/heatmap-hardened-skills-sae-forge.png}
  \caption*{\textbf{Spiral variant (plan-act)}: BUILD$\to$FIX and FIX$\to$BUILD
            cells are orange/red. The agent cycles. Worst run: 23 builds, \$5.67.}
\end{minipage}\hfill
\begin{minipage}[t]{0.48\textwidth}
  \includegraphics[width=\textwidth]{../figures/heatmap-hardened-skills-sae.png}
  \caption*{\textbf{Clean variant}: FIX$\to$BUILD cell is dimmer.
            Fewer cycles. Best run: 11 builds, \$3.49.}
\end{minipage}
\caption{The rework loop in the heatmap. Bright mutual probability between BUILD
         and FIX rows means the agent is cycling. The plan-act variant's two-phase structure
         produces this pattern when initial test writing fails.}
\end{figure}

%% ============================================================
\section{Step-Count Prediction: Does the Math Work?}
\label{sec:prediction}

The Markov framework is only useful if it predicts accurately.
We validated using leave-one-out cross-validation:
fit the transition matrix $P$ on $N-1$ runs per variant,
predict the held-out run's step count from the fundamental matrix $N = (I-Q)^{-1}$.

\begin{findingbox}
\textbf{Cross-validation result: zero mean bias, MAE of 8--20 steps.}

\begin{center}
\begin{tabular}{@{}lrrl@{}}
\toprule
\textbf{Variant} & \textbf{Actual/run} & \textbf{CV Predicted} & \textbf{MAE} \\
\midrule
\texttt{hardened}            & 103.0 & 103.0 & 19.7 \\
\texttt{hardened+kb}         &  83.3 &  83.3 &  8.3 \\
\texttt{hardened+sae}        &  80.3 &  80.3 & 13.3 \\
\texttt{hardened+skills}     & 101.7 & 101.7 & 11.0 \\
\texttt{hardened+skills+sae} &  75.0 &  75.0 & 10.3 \\
\bottomrule
\end{tabular}
\end{center}
\end{findingbox}

The predicted mean matches actual mean for all five variants with N=3 data.
The MAE of 8--20 steps represents run-to-run variance, not model error ---
the model captures the average, not the individual spiral.

\textbf{Why the prediction is only approximately Markov}:
A formal second-order test (9-state model, 1821 classified calls) found mean KL
divergence of 0.73 bits across 36 (previous, current) state pairs ---
first-order Markov is formally rejected ($p < 0.001$).
The violation is concentrated in recovery transitions:
FIX$\to$EXPLORE has 2.67 bits and VERIFY$\to$EXPLORE 3.37 bits of second-order information.
Knowing whether an agent arrived at FIX from its first build attempt or its
sixth dramatically changes what it explores next.

But over a 100-step trajectory, recovery transitions happen a handful of times
and the errors cancel. Hence: zero mean bias in cross-validation despite
technically non-Markov data.

%% ============================================================
\section{What This Experiment Cannot Show}
\label{sec:limits}

\begin{warningbox}
\textbf{The petclinic ceiling.}

Spring PetClinic is one of the most documented Spring Boot applications in existence.
All variants reach approximately 92--94\% instruction coverage.
Coverage is not a discriminator here --- the model knows this domain.

Consequences:
\begin{itemize}
  \item Knowledge injection (KB, skills) improves \emph{efficiency} (fewer steps) but
        cannot improve \emph{outcome quality} beyond the model's existing knowledge of petclinic.
  \item The T3 quality scores are meaningful --- they reflect how well-structured
        the tests are --- but they cannot reflect domain-knowledge improvement
        on a domain the model already knows.
  \item To prove that knowledge injection improves \emph{outcomes} on novel domains,
        we need a domain the model does not know well.
\end{itemize}
\end{warningbox}

\textbf{What petclinic does prove}:

\begin{itemize}
  \item Prompt hardening alone improves quality and reduces steps.
    This holds even when the model knows the domain perfectly.
  \item Structured execution (pre-analysis) reliably reduces steps on any codebase ---
    the pre-analysis benefit is architectural, not domain-specific.
  \item The JAR cluster is addressable via knowledge injection.
    Skills reduce it from 11\% to 1\% regardless of domain familiarity.
  \item The FIX/BUILD spiral risk is variant-specific.
    plan-act has it; \texttt{hardened+skills+sae} does not.
\end{itemize}

\textbf{What petclinic is}: the negative control.
It establishes the lower bound on knowledge injection's quality effect
(zero improvement on known domains) and the baseline behavioral signatures.

%% ============================================================
\section{What Comes Next: Fineract}
\label{sec:fineract}

The next experiment targets \textbf{Apache Fineract} (banking platform, Spring Boot 3.2.2,
Gradle 8.x). The \texttt{fineract-loan} module covers a bounded loan lifecycle domain ---
disbursement, repayment schedules, interest calculations, charge types.

This domain is not well-represented in model training data.
The model has seen petclinic; it has not seen Fineract's loan allocation rules.

\subsection*{What Fineract tests}

The central hypothesis: \textbf{on a novel domain, knowledge injection improves
test quality, not just test efficiency.}

The KB will contain 4 domain files:
\begin{itemize}
  \item \texttt{loan-lifecycle.md} --- disbursement $\to$ repayment $\to$ closure states
  \item \texttt{test-patterns.md} --- Spring Boot 3 test patterns, JPA, H2 in-memory
  \item \texttt{loan-charge-types.md} --- fee, penalty, and upfront charge classification
  \item \texttt{assertion-guide.md} --- how to write assertions for financial calculations
\end{itemize}

Fineract already has tests at roughly 40--65\% line coverage (module-dependent).
The task is to improve from that baseline, not write from scratch.
This is harder: the agent must understand what tests already exist
and write tests that extend coverage without duplicating them.

\subsection*{The experimental matrix}

6 variants $\times$ 3 runs = 18 sessions:
\begin{enumerate}
  \item \texttt{simple} --- baseline
  \item \texttt{hardened} --- structured prompt
  \item \texttt{hardened+kb} --- flat KB with loan domain files
  \item \texttt{hardened+skills} --- SkillsJars with same content
  \item \texttt{hardened+skills+sae} --- skills + pre-analysis
  \item \texttt{hardened+skills+sae+plan-act} --- two-phase
\end{enumerate}

\textbf{Success criterion}: at least one KB variant ($+$kb or $+$skills)
achieves $\geq$10\% T2-improvement (coverage gain above simple baseline).
If yes: knowledge injection proves outcome improvement on novel domains.
If no: efficiency is the only measurable benefit, and the thesis requires refinement.

%% ============================================================
\section{The Short Summary}

\begin{keybox}
\textbf{What the v2 experiment shows:}

\begin{enumerate}[leftmargin=1.5em,itemsep=4pt]
  \item \textbf{Prompt hardening is the single biggest quality driver.}
    $+$0.067 quality, $-$6\% steps. Free gain from structural discipline.

  \item \textbf{pre-analysis is efficient but shallow.}
    $-$22\% steps, quality back near baseline.
    Tight plans miss edge cases.

  \item \textbf{Skills fix pre-analysis's quality regression.}
    \texttt{hardened+skills+sae}: $-$31\% steps, quality matches \texttt{hardened}.
    Best tradeoff in the experiment.

  \item \textbf{KB and skills reduce JAR friction but not quality on known codebases.}
    On petclinic: efficiency gain, quality flat.
    On novel domains: untested.

  \item \textbf{plan-act is high variance.}
    Best quality ceiling, worst spiral risk.
    Its two-phase design amplifies both success and failure.

  \item \textbf{The Markov model predicts step counts accurately.}
    Zero mean bias in cross-validation despite formal rejection of first-order.
    Recovery transitions are non-Markov; averages are robust.
\end{enumerate}

\vspace{6pt}
\textbf{What comes next}: Apache Fineract --- a novel banking domain ---
to test whether the quality claim holds beyond a domain the model already knows.
\end{keybox}

\end{document}
